% !TEX root = ../main.tex
\section{Baseline comparison}
\label{sec:baselines}

The original work compares three methods at $N = 8$: standard-average FedNPG
\eqref{eq:stdfednpg}, FedNPG-ADMM, and first-order federated Proximal Policy Optimization (FedPPO) with clipping
parameter $0.2$. We implement all three in the same codebase, holding
everything except the server-side aggregation identical: same environments,
normalization, GAE, pooled critic and logged metrics. The comparison therefore
measures the algorithms, not the surrounding machinery.

\subsection{Standard-average FedNPG}

This arm solves \eqref{eq:stdfednpg} exactly by conjugate gradient on the
pooled Fisher, exploiting $\sum_i \bH_i v = N \cdot \bF_{\text{pooled}} v$.
The solve is matrix-free, but the \emph{metering} is not: the channel is
charged the full $d^2$ elements per agent per round.

The obvious objection is that a baseline meant to pay an $\Oh(d^2)$ cost ought
to pay it. Materializing $\sum_i \bH_i$ is impossible at Humanoid scale
(Table~\ref{tab:commcost}: \SI{2.1}{\peta\byte}), so a literal implementation
would not exist for the environment where the $\Oh(d^2)$ cost matters most.
The claim under test concerns \emph{uplink bytes}, not server FLOPs; the paper
measures ``the number of transmitted parameters with double precision in each
agent''. We therefore compute the identical direction cheaply and report the
true transmission cost. Dropping that metering would make FedNPG-ADMM appear
free by comparison, and the experiment would measure nothing.

\subsection{FedPPO}

The paper names FedPPO and gives its clipping parameter, and specifies nothing further, neither the aggregation scheme nor the local
epoch count. We implement the
FedAvg reading: each agent runs local PPO epochs from the broadcast $\bthe$
and uplinks its parameter delta, and the server averages. The alternative, a
single clipped-surrogate gradient step per round, would parallel
Algorithm~\ref{alg:fednpgadmm}'s structure more neatly, but it strips PPO of
the multi-epoch minibatch reuse that its clipping exists to make safe. Beating
that version would establish nothing. The choice is recorded in every run's
metadata.

\subsection{Communication cost, measured}

\begin{table}[t]
  \centering
  \caption{Measured per-iteration uplink, Swimmer-v4 with four agents
  ($d = \num{4868}$), at 64 bits per parameter. Measured values match the closed
  form exactly.}
  \label{tab:measuredcomm}
  \begin{tabular}{llrr}
    \toprule
    Method & Transmitted & Bytes/iteration & Relative \\
    \midrule
    Standard-average FedNPG & $\bH_i$ ($d^2$) $+\ \bg_i$ & \num{758473344} & $2434\times$ \\
    FedNPG-ADMM             & $\by_i,\ \bg_i$ (2 vectors) & \num{311552}   & $2\times$ \\
    FedPPO                  & $\Delta\bthe_i$ (1 vector)  & \num{155776}   & $1\times$ \\
    \bottomrule
  \end{tabular}
\end{table}

One consequence of Table~\ref{tab:measuredcomm} deserves emphasis, because the
paper's framing invites a reader to assume otherwise: FedNPG-ADMM is not the
cheapest of the three methods. It is $\Oh(d)$, but with twice the constant of
FedPPO, since it transmits two vectors per round where FedPPO transmits one.
The argument for the method therefore rests on what it achieves per round
rather than on what it costs per round, and a figure plotted against iterations
cannot show that. We consequently plot reward against \emph{cumulative bytes},
which is the axis on which the two effects can be read together.

\TODO{insert Figure 3 analogue -- reward vs iterations and reward vs cumulative
bytes, three methods, N=8, three environments, five seeds -- once the 30
baseline runs complete}

\TODO{table of final rewards for the three methods at N=8}
